\documentclass[11pt,a4paper]{article}

% --- Packages ---
\usepackage[utf8]{inputenc}
\usepackage[T1]{fontenc}
\usepackage{lmodern}
\usepackage[margin=25mm]{geometry}
\usepackage{amsmath,amssymb,amsthm,amsfonts}
\usepackage{mathtools}
\usepackage{mathrsfs}
\usepackage{booktabs}
\usepackage{tabularx}
\usepackage{enumitem}
\usepackage{hyperref}
\usepackage{xcolor}

% --- Hyperref Setup ---
\hypersetup{
    colorlinks=true,
    linkcolor=blue!60!black,
    citecolor=blue!60!black,
    urlcolor=blue!70!black,
    pdftitle={The Form of Being (Ordo Forma Est): Holographic Structure of Spacetime, Emergence, and Consciousness, and the Ontological Necessity of Mathematical Form},
    pdfauthor={H. Ota}
}

% --- Theorem Environments ---
\theoremstyle{definition}
\newtheorem{definition}{Definition}[section]
\newtheorem{axiom}{Axiom}[section]
\newtheorem{remark}{Remark}[section]

\theoremstyle{plain}
\newtheorem{theorem}{Theorem}[section]
\newtheorem{proposition}{Proposition}[section]
\newtheorem{corollary}{Corollary}[section]

% --- Title & Author ---
\title{\textbf{The Form of Being (\textit{Ordo Forma Est}):}\\
\Large Holographic Structure of Spacetime, Emergence, and Consciousness,\\
and the Ontological Necessity of Mathematical Form}
\author{Hiroshi Ota\thanks{Homepage: \url{https://www.kae-made.jp}, Research Notes: \url{https://note.com/kae_made/membership}}\\
\small Independent Researcher / Knowledge \& Experience\\
\small \texttt{master@kae-made.jp}}
\date{November 25, 2026}

\begin{document}

\maketitle

\begin{abstract}
In the preceding four papers (Paper 0, Paper 1, Paper 2, and Paper 3), the author formulated a relation-based ontology grounded in spans rather than isolated objects (``things'') and a universe composed of multiple ``Fields of Sense.'' Across these works, it was proved that: (i) static semantic structures within a single field of sense form a conceptual information model represented as a quotient 2-category $\mathcal{C} \equiv CG/\mathcal{R}$; (ii) dynamic behaviors are governed by state models, a symmetric monoidal category $\mathcal{DF}$ of acyclic data-flow action models, and a universal execution semantics $\mathcal{E}_{\mathrm{exec}}$ defined strictly over references to form ($\mathcal{C}_{\mathrm{FORM}}$); (iii) parallel fields of sense are connected via ``bridges as fields of sense''; and (iv) unique structural transformations under boundary conditions provided by ``aspectual fields of sense'' are \textbf{completely isomorphic to propagation operators via Green's functions} in mathematical physics.

As the final installment of this five-part series, the present paper (Paper 4) integrates all these mathematical and engineering apparatuses to prove that the formalism of conceptual modeling is not merely an epistemological or descriptive tool invented by humans to represent the world, \textbf{but rather the universal ontological form by which this world itself comes into being and manifests (\textit{Ordo Forma Est}: the form of being is the model)}. Specifically, this paper establishes four ultimate propositions. First, under the plenitude of forms satisfying the category of form $\mathcal{C}_{\mathrm{FORM}}$, structures containing contradictions or ambiguities immediately self-eliminate into the initial object $0$ (the empty space) of the internal logic of a topos, while only consistent trajectories autonomously evolve (self-unfold) over time under $\mathcal{E}_{\mathrm{exec}}$ (\textbf{Theorem of Self-Elimination of Form}). Second, the emergence of continuous higher-dimensional bulk spacetime and physical laws from discrete ``relations and forms (\textit{Ordo Forma Est})'' inscribed on a lower-dimensional boundary screen is \textbf{categorically completely isomorphic to the Holographic Principle} of modern physics (\textbf{Holographic Isomorphism Theorem}). Third, ``phase transitions'' are formulated as horizontal field-switching triggered by boundary thresholds of hidden bridges (the replacement of water instances with ice instances), while ``emergence'' is formulated as a phenomenon in which a pullback via a \textbf{hidden real bridge ($\mathcal{B}^*$)} connecting a source field of sense and an aspectual field of sense \textbf{automatically and continuously triggers a unique structural transformation into a higher-level target field of sense without any observer (homunculus)}---thereby unifying the emergent reality of life, consciousness, thought, and even fictional or narrative worlds. Fourth, given that the human ``field of sense of consciousness and thought'' emerges from within a holographically emergent world, the fact that the observed world obeys mathematical principles is proved to be an \textbf{ontological necessity (a self-evident truism)}, while a strict suspension of judgment (\textit{epoch\'e}) is maintained regarding the origin of the foundational ``boundary screen (field of sense of qubit information)'' that serves as the matrix of Einsteinian spacetime, as it lies in a realm beyond verification or falsification.
\end{abstract}
\section{Introduction}

\subsection{Problem Statement: From a ``Tool of Description'' to the ``Form of Being Itself (\textit{Ordo Forma Est})''}
Humanity has long regarded conceptual modeling and mathematical modeling as tools for human beings to observe, abstract, and describe an external ``collection of things'' presumed to exist prior to description. However, what became clear through the preceding four papers (Papers 0--3) is a startling fact diametrically opposed to this conventional wisdom.

Specifically, the structures at which human beings inevitably arrive when stripping away ambiguity to the extreme in order to consistently describe real-world business systems or natural phenomena---namely, the conceptual information model as a quotient 2-category (Paper 1), state models marking proper time and local inertial frames alongside the symmetric monoidal category $\mathcal{DF}$ of acyclic string diagrams, the universal execution semantics $\mathcal{E}_{\mathrm{exec}}$ over the category of form $\mathcal{C}_{\mathrm{FORM}}$ (Paper 2), and unique structural transformations via aspectual fields of sense (Paper 3)---turned out to be \textbf{not merely metaphorical analogies, but 100\% categorically isomorphic to ``field theories in physics (configuration spaces, phase spaces, equations of motion, and principles of time evolution)'' and ``Green's function propagation kernels in mathematical physics.''}

Why does the ``form of conceptual modeling,'' which supposedly originated in human software design and systems analysis, coincide so completely with the mathematical structures of natural physical laws and boundary value problems? There can be only one answer. As the author concludes, \textbf{it is not that human beings impose form onto the world; rather, this world (being itself) is constituted from the very beginning by ``relations and form (\textit{Ordo Forma Est}).''}

\subsection{Four Ultimate Questions of This Paper}
As the concluding volume of the five-part series, this paper (Paper 4) provides rigorous mathematical answers to the following four ultimate ontological and scientific questions:

\begin{enumerate}
    \item \textbf{The Question of Cosmic Emergence via Plenitude and Self-Elimination of Form:} When all schemas and states satisfying the form of conceptual modeling ($\mathcal{C}_{\mathrm{FORM}}$) are potentially permitted, can the mechanism by which inconsistent structures self-collapse into the initial object $0$, causing consistent cosmic histories and physical laws to autonomously self-unfold, be mathematically proved?
    \item \textbf{The Question of Structural Isomorphism with the Holographic Principle:} Is the structure whereby continuous higher-dimensional bulk spacetime and gravity emerge from discrete ``relations and forms'' inscribed on a boundary screen categorically completely isomorphic to the Holographic Principle of modern physics (such as the AdS/CFT correspondence and \textit{It from Bit})?
    \item \textbf{The Question of the True Nature of Phase Transition, Emergence, and Consciousness:} Can ``emergence'' and ``consciousness,'' which traditional science and philosophy have mystified, be completely elucidated as phenomena in which a pullback via a hidden real bridge ($\mathcal{B}^*$) and a unique structural transformation into a target field of sense \textbf{automatically and continuously trigger without any observer (homunculus)}? Furthermore, from this perspective, how is the reality of fictional worlds such as novels, films, and manga positioned?
    \item \textbf{The Question of the Final Answer to ``Why Do Mathematical Principles Hold?'' and the Suspension of Judgment:} While providing an irrefutable final answer to what Eugene Wigner called ``the unreasonable effectiveness of mathematics in the natural sciences'' from the structure wherein consciousness emerges from within a holographically emergent world, where should the scientific and epistemological boundary line be drawn regarding the origin of the foundational ``boundary screen (field of sense of qubit information)''?
\end{enumerate}

\subsection{Structure of This Paper}
This paper is organized into seven sections:

\begin{itemize}
    \item \textbf{Section 2 (Unification of the Tetralogy and the Fundamental Principle \textit{Ordo Forma Est}):} Formulates all concepts from Papers 0--3 as a unified category-theoretic architecture $\mathfrak{U}_{\mathrm{OFE}}$, and proves the Principle of Plenitude of Forms and the Theorem of Self-Elimination via the initial object $0$.
    \item \textbf{Section 3 (Structural Isomorphism Theorem with the Holographic Principle of Modern Physics):} Defines a category-theoretic specification of the Holographic Principle and proves a complete isomorphism theorem (along with a one-to-one correspondence table) between it and the emergence of continuous bulk spacetime from discrete boundary forms.
    \item \textbf{Section 4 (The True Nature of Phase Transition, Emergence, and Consciousness---Hidden Bridges and Automatically Triggered Transformations):} Presents a category-theoretic analysis of the ``water-to-ice phase transition,'' formulates ``emergence'' as a hidden real bridge $\mathcal{B}^*$ coupled with an automatically triggered unique transformation, critiques conventional debates on emergence, examines ``consciousness'' as an automatically triggered transformation without an observer, and discusses the ontological status of narrative worlds and fictional characters (such as Natsuhiko Kyogoku's \textit{T\=ofu-koz\=o}).
    \item \textbf{Section 5 (Why Do Mathematical Principles Hold?---Ontological Necessity and Suspension of Judgment):} Provides a final answer to Wigner's enigma via the Internal Observer Theorem, demonstrates the mathematical completion of the philosophies of Immanuel Kant and Kitaro Nishida, and explicitly states a suspension of judgment (\textit{epoch\'e}) regarding the origin of the primordial ``field of sense of qubit information (boundary screen).''
    \item \textbf{Section 6 (Responses to Anticipated Objections and Practical/Scientific Implications):} Offers mathematical responses to five anticipated philosophical and scientific objections alongside an open call for verification as a working paper, and discusses implications for quantum gravity theory, the conditions for true artificial consciousness (establishing a decisive boundary with current LLMs), and a unified language bridging the humanities and natural sciences.
    \item \textbf{Section 7 (Summary of the Five-Part Series (Papers 0--4) and Conclusion):} Summarizes the theoretical achievements of the entire series and concludes with the declaration of \textit{Ordo Forma Est}.
\end{itemize}
\section{Unification of the Tetralogy and the Fundamental Principle \textit{Ordo Forma Est}}

\subsection{Formulation of the Unified Category-Theoretic Architecture $\mathfrak{U}_{\mathrm{OFE}}$}
First, all mathematical components established from Paper 0 through Paper 3 are formulated into a single unified category-theoretic architecture.

\begin{definition}[Unified Architecture of the Form of Being $\mathfrak{U}_{\mathrm{OFE}}$]
The formal system $\mathfrak{U}_{\mathrm{OFE}}$ of being and phenomena in the universe $\mathcal{U}$ is defined as a 5-tuple of structures:
\[
\mathfrak{U}_{\mathrm{OFE}} \equiv \left(\mathcal{C}_{\mathrm{FORM}},\; \mathcal{C}_{\mathrm{SUB}},\; \mathfrak{B},\; \mathcal{E}_{\mathrm{exec}},\; \mathfrak{T}\right).
\]
\begin{enumerate}
    \item \textbf{Category of Form ($\mathcal{C}_{\mathrm{FORM}}$):} The universal reference category whose objects and morphisms comprise the structural framework shared by all fields of sense---namely, concept classes, 7 primitive data types, identifier monomorphisms, 4 types of multiplicity, mutually exclusive direct-sum \texttt{Is-a} hierarchies, span associative classes, state transition tables, and DAG wiring of 16 elementary process components (Paper 1; Paper 2, Definition 4.1; Paper 3, Definition 3.1).
    \item \textbf{Category of Substantive Fields of Sense ($\mathcal{C}_{\mathrm{SUB}}$):} The category of individual substantive fields of sense $\mathcal{M}$, each possessing an instance functor $\Phi_{\mathcal{M}} \in \mathrm{Valid}(\mathcal{C}_{\mathrm{FORM}})$ over the category of form $\mathcal{C}_{\mathrm{FORM}}$ as its schema, equipped with a quotient 2-category $\mathcal{C} \equiv CG/\mathcal{R}$ and a symmetric monoidal category $\mathcal{DF}$.
    \item \textbf{Bridges and Aspectual Fields ($\mathfrak{B}$):} The system comprising ``bridges as fields of sense $\mathcal{B}$'' that connect parallel fields of sense via purely structural correspondence morphisms (functors) devoid of semantic phrases and higher-order operators (\texttt{=>}, \texttt{<=>}, \texttt{REQUIRE}) (Paper 3, Section 2), together with ``aspectual fields of sense $\mathcal{M}_{\mathrm{Aspect}}$'' that treat form itself ($\mathcal{C}_{\mathrm{FORM}}$) as individuals and impose boundary conditions as fiber products ($\times_{\mathcal{C}_{\mathrm{FORM}}}$) (Paper 3, Section 3).
    \item \textbf{Universal Execution Semantics ($\mathcal{E}_{\mathrm{exec}}$):} The state transition functor uniquely and deterministically defined solely over the category of form $\mathcal{C}_{\mathrm{FORM}}$ without inspecting any domain-specific content---incorporating Run-to-Completion, self-event priority, FIFO ordering with parallel commutativity, and aspectual pullbacks (Paper 2, Theorem 4.1).
    \item \textbf{Structural Transformation System ($\mathfrak{T}$):} The system of unique transformation functors $\mathcal{T}_{\mathcal{M}_{\mathrm{Aspect}}}$, isomorphic to Green's function propagation kernels, induced from the pullback $\mathcal{M}_{\mathrm{Src}} \times_{\mathcal{C}_{\mathrm{FORM}}} \mathcal{M}_{\mathrm{Aspect}}$ to map source forms unidirectionally into instance states satisfying a target schema $\mathcal{C}_{\mathrm{Targ}}$ (Paper 3, Theorems 4.1 and 5.5).
\end{enumerate}
\end{definition}

\noindent\textbf{Fundamental Principle (\textit{Ordo Forma Est}: The Form of Being Is the Model):}
The unified architecture $\mathfrak{U}_{\mathrm{OFE}}$ does not enter into any specific content (\textit{Matter}), such as what concrete substances, organisms, or social institutions exist in this world. However, for any entity whatsoever, the \textbf{conditions of possibility (the form of being: \textit{Forma})} required for it to exist consistently, undergo causal change, and manifest in coordination with other domains are exhaustively and precisely determined by $\mathfrak{U}_{\mathrm{OFE}}$.

\subsection{Principle of Plenitude of Forms and Theorem of Self-Elimination via the Initial Object $0$}
Next, the most fundamental dynamics governing how concrete cosmic histories and complex structures autonomously arise under this formal system $\mathfrak{U}_{\mathrm{OFE}}$ is formulated.

In Paper 1 (Theorem 5.2) and Paper 2 (Theorem 5.2), it was proved that for any consistent schema category $\mathcal{C}$, a non-trivial valid state $\mathcal{F} \in \mathrm{Valid}(\mathcal{C})$ and normal operational trajectories constructively exist, whereas structures containing multiplicity violations, out-of-range values, division by zero, \texttt{Can't Happen} transitions, or bilateral schema inconsistencies collapse without exception into the \textbf{initial object $0$ (empty space / contradictory state)} of the internal logic of a topos.

\begin{axiom}[Principle of Plenitude of Forms]
All schemas and states that consistently satisfy the axiom system of the category of form $\mathcal{C}_{\mathrm{FORM}}$, along with all bridges and boundary-condition aspects connecting them, are potentially permitted as mathematical realities (plenitude of forms).
\end{axiom}

\begin{theorem}[Theorem of Self-Elimination of Form and Autonomous Self-Unfolding of the Universe]
Under the Principle of Plenitude of Forms (Axiom 2.1) and the universal execution semantics $\mathcal{E}_{\mathrm{exec}}$, the following autonomous dynamics inevitably arise without assuming any transcendent designer or external manual intervention (an observer pressing a button):
\begin{enumerate}
    \item \textbf{Immediate Self-Elimination of Inconsistent Structures (Collapse into the Initial Object $0$):} Any form or trajectory containing internal contradictions, structural breakdown due to linguistic polysemy (Paper 0, Theorem 3.2), arrival at singularities (Paper 2, Corollary 4.1), or violations of bilateral schema conservation (Paper 3, Axiom 2.1) immediately self-collapses into the initial object $0$ (empty space) of the internal logic of a topos and cannot persist as a phenomenon.
    \item \textbf{Autonomous Self-Unfolding of Consistent Trajectories (Automatic Driving of Equational Logic):} Consistent fields of sense and normal operational trajectories that escape collapse into the initial object $0$ autonomously and continuously evolve (self-unfold) over time without external intervention, governed by the equational logic of $\mathcal{E}_{\mathrm{exec}}$ (automatic firing of actions upon satisfaction of input conditions and chaining of self-events).
\end{enumerate}
\end{theorem}

\begin{proof}[Proof Sketch]
By Paper 2 (Theorems 3.1 and 4.1), the data-flow action model is not an imperative procedure but an equational logic triggered by resource satisfaction in linear logic, and $\mathcal{E}_{\mathrm{exec}}$ induces a unique deterministic transition system $\mathrm{Valid}(\mathcal{C}) \Rightarrow \mathrm{Valid}(\mathcal{C}) \sqcup \{0\}$ over the category of form $\mathcal{C}_{\mathrm{FORM}}$. In the space of forms granted by Axiom 2.1, regions failing to satisfy consistency constraints map to $0$ (the empty set) via pullback failure and vanish, while only consistent regions satisfying $\mathrm{Valid}(\mathcal{C}) \neq \emptyset$ repeatedly undergo closed self-transitions as normal operational trajectories $\mathcal{F}_0 \to \mathcal{F}_1 \to \cdots$. Consequently, the self-elimination of contradictions and the autonomous self-unfolding of consistent trajectories follow as a mathematical necessity.
\end{proof}

\section{Structural Isomorphism Theorem with the Holographic Principle of Modern Physics}

\subsection{Category-Theoretic Specification of the Holographic Principle $\mathcal{G}_{\mathrm{Holo}}$}
The unified architecture $\mathfrak{U}_{\mathrm{OFE}}$ formulated in the preceding section is fundamentally a formal system composed of discrete ``relations (spans),'' ``identifiers,'' ``state transitions,'' and ``elementary process components.'' How, then, do the ``smooth, continuous three-dimensional space and time (Einsteinian spacetime)'' and the ``continuous dynamics of gravity and fields'' experienced in everyday life and in physics relate to this discrete formal system?

The decisive key to answering this question lies in the \textbf{Holographic Principle} of modern theoretical physics (quantum gravity theory, black hole thermodynamics, and the AdS/CFT correspondence pioneered by Maldacena et al.). According to the Holographic Principle, all information regarding continuous physical phenomena---including gravity---occurring within a $(d+1)$-dimensional spatial region (\textbf{the Bulk}) is completely encoded by discrete degrees of freedom (qubit information and entanglement structures) inscribed on a lower-dimensional ($d$-dimensional) \textbf{Boundary Screen} enclosing that region, such that the spacetime geometry of the higher-dimensional bulk itself emerges from discrete relational structures on the boundary (\textbf{Emergent Spacetime}).

John Archibald Wheeler once encapsulated this information-ontological insight in the motto ``\textit{It from Bit}'' (that all matter and spacetime arise from binary information). From the perspective of the present framework, however, neither meaning nor space can arise from mere unstructured strings of bits (0s and 1s). What is inscribed upon the boundary screen is nothing other than \textbf{Form equipped with identifiers and relations (spans)} (\textit{Ordo Forma Est}: \textit{It from Form / Relation}).

\begin{definition}[Category-Theoretic Specification of the Holographic Principle $\mathcal{G}_{\mathrm{Holo}}$]
The formal structure $\mathcal{G}_{\mathrm{Holo}}$ of bulk spacetime emergence via the Holographic Principle in modern physics is categorically specified as the following 4-tuple:
\[
\mathcal{G}_{\mathrm{Holo}} \equiv \left(\partial\mathcal{M}_{\mathrm{Boundary}},\; \mathcal{E}_{\mathrm{Entangle}},\; \mathcal{K}_{\mathrm{Bulk\text{-}Boundary}},\; \mathcal{M}_{\mathrm{Bulk}}\right).
\]
\begin{enumerate}
    \item \textbf{Discrete Degrees of Freedom on the Boundary Screen ($\partial\mathcal{M}_{\mathrm{Boundary}}$):} Discrete quantum states and degrees of freedom arranged on a lower-dimensional boundary, together with their local correlation structures.
    \item \textbf{Boundary Time Evolution and Entanglement Structure ($\mathcal{E}_{\mathrm{Entangle}}$):} Informational coupling (entanglement) among boundary degrees of freedom and unitary time-evolution rules on the boundary.
    \item \textbf{Bulk-to-Boundary Propagator / Reconstruction Kernel ($\mathcal{K}_{\mathrm{Bulk\text{-}Boundary}}$):} A Green's-function-type propagation kernel (such as the HKLL reconstruction kernel) that reconstructs local fields and metrics inside the higher-dimensional bulk from local states on the boundary.
    \item \textbf{Emergent Higher-Dimensional Continuous Bulk Spacetime ($\mathcal{M}_{\mathrm{Bulk}}$):} A higher-dimensional spacetime manifold equipped with a smooth metric structure, metric tensor, and gravity, emerging as the coarse-grained continuous limit of informational correlations (entanglement) on the boundary.
\end{enumerate}
\end{definition}

\subsection{Holographic Isomorphism Theorem and Comparison Table}
By integrating Paper 2 (Theorem 5.5: construction of continuous dynamics via projective limits and sheaves) and Paper 3 (Theorem 5.5: Green's Function Isomorphism Theorem), it follows that the present system $\mathfrak{U}_{\mathrm{OFE}}$ is mathematically completely isomorphic to the Holographic Principle $\mathcal{G}_{\mathrm{Holo}}$.

\begin{theorem}[Holographic Isomorphism Theorem]
Within the unified architecture $\mathfrak{U}_{\mathrm{OFE}}$, let $\mathfrak{B}_{\mathrm{OFE}}$ denote the continuous spacetime and physical bulk system that emerges by:
\begin{enumerate}
    \item taking as its foundation the discrete ``relations (spans) and forms ($\mathcal{C}_{\mathrm{FORM}}$)'' inscribed on a boundary screen together with discrete state transitions governed by the universal execution semantics $\mathcal{E}_{\mathrm{exec}}$, and
    \item applying the projective limit ($\varprojlim_{\varepsilon \to 0}$) over boundary-condition aspects $\mathcal{V}_{\varepsilon}$ and sheaf gluing (Paper 2, Theorem 5.5), along with coend reconstruction operators via Green's function propagation kernels (Paper 3, Theorem 5.5).
\end{enumerate}
Then $\mathfrak{B}_{\mathrm{OFE}}$ is \textbf{categorically equivalent (structurally completely isomorphic: $\mathfrak{B}_{\mathrm{OFE}} \simeq \mathcal{G}_{\mathrm{Holo}}$)} to the category-theoretic specification $\mathcal{G}_{\mathrm{Holo}}$ of the Holographic Principle in modern physics.
\end{theorem}

\begin{table}[htbp]
\centering
\small
\renewcommand{\arraystretch}{1.35}
\begin{tabularx}{\textwidth}{@{} >{\raggedright\arraybackslash}p{0.31\textwidth} >{\raggedright\arraybackslash}p{0.31\textwidth} >{\raggedright\arraybackslash}X @{}}
\toprule
\textbf{Form of Being System $\mathfrak{B}_{\mathrm{OFE}}$ (This Series)} & \textbf{Holographic Principle of Modern Physics $\mathcal{G}_{\mathrm{Holo}}$} & \textbf{Mathematical \& Ontological Isomorphic Correspondence} \\
\midrule
\textbf{1. Forms and Relations on the Boundary Screen}\newline (Quotient 2-category $\mathcal{C} \equiv CG/\mathcal{R}$, spans, qubit-like discrete states) & \textbf{Lower-Dimensional Boundary Degrees of Freedom $\partial\mathcal{M}_{\mathrm{Boundary}}$}\newline (Boundary qubits, local operators of conformal field theory) & Bulk spacetime is not presupposed \textit{a priori}; only discrete distinctions (identifiers) and relations (spans) inscribed on the boundary constitute primary reality (\textit{Ordo Forma Est}). \\
\midrule
\textbf{2. Universal Execution Semantics and Bridge Coupling}\newline ($\mathcal{E}_{\mathrm{exec}}$, symmetric monoidal category $\mathcal{DF}$, \texttt{<=>} synchronization) & \textbf{Boundary Time Evolution and Entanglement $\mathcal{E}_{\mathrm{Entangle}}$}\newline (Unitary evolution and quantum entanglement structure on the boundary) & Causal chains and forced synchronization (entanglement) among discrete boundary elements generate the prototypes of spatial adjacency and distance (topology). \\
\midrule
\textbf{3. Green's Function Transformation Kernel and Projective Limit}\newline (Coend integral $\int^{\xi} K(x,\xi) \times \Phi(\xi)$ and $\varprojlim_{\varepsilon \to 0}$) & \textbf{Bulk-to-Boundary Propagator $\mathcal{K}_{\mathrm{Bulk\text{-}Boundary}}$}\newline (Green's function kernel reconstructing bulk fields from boundary states) & Local values inside the bulk are uniquely reconstructed from discrete boundary forms via limit operations on boundary-condition aspects and superpositions of Green's function kernels. \\
\midrule
\textbf{4. Emergent Continuous Spacetime and Physical World}\newline (Sheaf gluing $\mathcal{F}_{\mathrm{Sheaf}}$ and physical theory isomorphism $\mathcal{P}_{\mathrm{Phys}}$) & \textbf{Emergent Higher-Dimensional Continuous Bulk $\mathcal{M}_{\mathrm{Bulk}}$}\newline (Emergence of Einsteinian spacetime, metric, and gravitational fields) & Smooth continuous spacetime and physical laws (Paper 2, Theorem 5.6) emerge without information loss as the limit structure of discrete relations and forms. \\
\bottomrule
\end{tabularx}
\caption{Isomorphic Correspondence between the Form of Being System $\mathfrak{B}_{\mathrm{OFE}}$ and the Holographic Principle $\mathcal{G}_{\mathrm{Holo}}$}
\end{table}

\begin{proof}[Proof Sketch]
As shown in Paper 2 (Remark 1.1), the foundation of the pure conceptual modeling system $\mathfrak{U}_{\mathrm{OFE}}$ contains no \textit{a priori} continuous spacetime; rather, there exist only discrete sets of individuals, relational spans, and proper-time event queues on fields of sense corresponding to the boundary. The couplings among these discrete elements (parallel composition $\otimes$ and \texttt{<=>} synchronization corresponding to bridges and entanglement) define topological connectivity, while the coend integral formula of Paper 3 (Theorem 5.5), $\mathcal{M}_{\mathrm{Targ}}^*(x) \cong \int^{\xi \in \mathcal{C}_{\mathrm{FORM}}} K_{\mathcal{M}_{\mathrm{Aspect}}}(x, \xi) \times \Phi_{\mathrm{Src}}(\xi)$, reconstructs responses at interior coordinates $x$ from discrete forms on the boundary. Furthermore, by Paper 2 (Theorem 5.5), under the projective limit $\varprojlim_{\varepsilon \to 0} (\mathcal{M} \times_{\mathcal{C}_{\mathrm{FORM}}} \mathcal{V}_{\varepsilon})$ as the resolution parameter $\varepsilon \to 0$ and the sheaf gluing conditions, this reconstructed state-evolution system becomes categorically equivalent to a physical field theory $\mathcal{P}_{\mathrm{Phys}}$ on a smooth continuous manifold (Paper 2, Theorem 5.6). This is completely isomorphic to the mathematical structure of the Holographic Principle $\mathcal{G}_{\mathrm{Holo}}$, in which continuous bulk spacetime emerges from discrete quantum information and entanglement on the boundary via Green's function propagation kernels.
\end{proof}

\begin{corollary}[Sublimation from \textit{It from Bit} to \textit{It from Form / Relation}]
The continuous spacetime and gravity described by Einstein's general theory of relativity are not fundamental substances of the universe, \textbf{but secondary (emergent) bulk structures that emerge holographically from ``discrete relations and forms (\textit{Ordo Forma Est}) inscribed on a boundary screen.''}
\end{corollary}

\section{The True Nature of Phase Transition, Emergence, and Consciousness---Hidden Bridges and Automatically Triggered Transformations}

\subsection{Category-Theoretic Mechanism of Phase Transition: Reinterpreting ``Water and Ice''}
In the preceding section, it was demonstrated how continuous bulk spacetime arises holographically from discrete boundary forms. How, then, are qualitative changes occurring within that world---namely, \textbf{``Phase Transition''} and \textbf{``Emergence''}---distinguished and unified as mathematical mechanisms within the present framework $\mathfrak{U}_{\mathrm{OFE}}$?

First, consider the most familiar phase transition in the material world: \textbf{the transition from liquid water to solid ice}. In the commonsense view of matter (thing-based ontology), this is interpreted as: ``A single substance (`thing') called `water ($\mathrm{H_2O}$)' exists first, and upon cooling, it changes its form from the property of liquidity to the property of solidity.'' From the category-theoretic perspective of the present system (Papers 0--3), however, a completely different and rigorous structure comes into view:

\begin{enumerate}
    \item \textbf{Two Distinct Fields of Sense Governed by Different Schema Categories:}
    The ``field of sense of liquids ($\mathcal{M}_{\mathrm{Liquid}}$ / schema category $\mathcal{C}_{\mathrm{Liquid}}$),'' which is governed by constraints (a constitution) such as fluidity, viscosity, and volume conservation, and the ``field of sense of solids ($\mathcal{M}_{\mathrm{Solid}}$ / schema category $\mathcal{C}_{\mathrm{Solid}}$),'' which is governed by constraints such as crystal lattice structure, rigidity, and resistance to shear stress, are \textbf{distinct fields of sense governed by mutually different schema categories}.
    \item \textbf{Thresholds of Hidden Bridges and Collapse into the Initial Object $0$ (Horizontal Field-Switching):}
    Behind the scenes of the material world, a \textbf{hidden real bridge ($\mathcal{B}^*$)}---invisible directly to the human eye---spans between these fields of sense and the ``field of sense of thermodynamics and temperature ($\mathcal{M}_{\mathrm{Thermo}}$),'' imposing boundary invariant conditions (\texttt{REQUIRE}: e.g., $T > 0^\circ\mathrm{C}$). The instant the ambient temperature drops below $0^\circ\mathrm{C}$, the constraints on the $\mathcal{M}_{\mathrm{Liquid}}$ side break down, and \textbf{the existence of that water instance in the liquid field of sense suddenly collapses (vanishes) into the initial object $0$ (empty space) of the internal logic of a topos}. Simultaneously, on the opposite side of the same hidden bridge and the boundary-condition aspects of mass-energy conservation, \textbf{a new ``ice instance'' suddenly appears (manifests) via a creation event (\texttt{CREATE}) within the solid field of sense $\mathcal{M}_{\mathrm{Solid}}$}.
    \item \textbf{Why Human Beings Perceive It as the ``Same Water'':}
    What human beings directly perceive in everyday sensation and observation is solely \textbf{``phenomena as intensions (coldness, transparency, occupation of the same spatial coordinates, continuity of mass, etc.)''} projected from each field of sense into the field of consciousness via domain functions and sensory organs. Moreover, within the human ``field of sense of everyday language and common sense ($\mathcal{M}_{\mathrm{Everyday}}$),'' there exists a schema that bundles together groups of phenomena exhibiting continuous mass and spatial location as a single individual labeled ``water.'' Consequently, even though a discontinuous field-switching---the collapse and disappearance of a liquid instance alongside the novel emergence of an ice instance in a separate domain---is occurring in the underlying reality, human observers who perceive only phenomena as intensions interpret it as ``the same water merely changing its state.''
\end{enumerate}

\subsection{Category-Theoretic Formulation of Emergence: Hidden Real Bridges and Automatically Triggered Structural Transformations}
Whereas the phase transition described above is a \textbf{horizontal transition between fields}, in which an instance in one field collapses into the initial object $0$ and switches to an instance in another field, \textbf{``Emergence''}---known as the appearance of life phenomena, consciousness and thought, or complex social orders---is a \textbf{vertical structural transformation}, in which the lower-level source field of sense persists without collapsing while its ``form ($\mathcal{C}_{\mathrm{FORM}}$)'' is dimensionally converted and simultaneously arises as a ``state (instance)'' of a novel, higher-level field of sense.

In conventional science and philosophy, emergence has often been spoken of as ``the mysterious appearance of a wholeness exceeding the sum of its parts.'' Viewed from the formal system established in Paper 3 (Sections 3 and 4), however, the mechanism of emergence is a purely mathematical consequence devoid of any mystery whatsoever.

\begin{theorem}[Mechanism of Emergence via Hidden Real Bridges]
``Emergence'' from a lower-level substantive field of sense (the source $\mathcal{M}_{\mathrm{Src}}$: e.g., a field of physicochemical reactions or neural firings) to a novel higher-level field of sense (the target $\mathcal{M}_{\mathrm{Targ}}$: e.g., a field of sense of life or of consciousness and thought) is a phenomenon that inevitably occurs whenever the following three structures are in place:
\begin{enumerate}
    \item A unique structural transformation functor $\mathcal{T}_{\mathcal{M}_{\mathrm{Aspect}}}$ is defined between an \textbf{aspectual field of sense $\mathcal{M}_{\mathrm{Aspect}}$ (boundary condition)} referring to the source form ($\mathcal{C}_{\mathrm{FORM}}$) and the target field of sense $\mathcal{M}_{\mathrm{Targ}}$ (Paper 3, Theorem 4.1).
    \item Between the source $\mathcal{M}_{\mathrm{Src}}$ (referred to via the aspectual field of sense) and $\mathcal{M}_{\mathrm{Aspect}}$, a fiber product (pullback $\mathcal{M}_{\mathrm{Src}} \times_{\mathcal{C}_{\mathrm{FORM}}} \mathcal{M}_{\mathrm{Aspect}}$) holds via a \textbf{``Hidden Real Bridge ($\mathcal{B}^*$)'' that already exists in reality as part of the world's order (merely not yet recognized or verbalized by human beings)}.
    \item The instant the above pullback condition is satisfied, under the equational logic of the universal execution semantics $\mathcal{E}_{\mathrm{exec}}$, \textbf{the unique structural transformation $\mathcal{T}_{\mathcal{M}_{\mathrm{Aspect}}}$ automatically triggers (without anyone needing to press a switch)}, autonomously bringing forth the instance state of the target field of sense $\mathcal{M}_{\mathrm{Targ}}$.
\end{enumerate}
\end{theorem}

\begin{remark}[Epistemological Limitations of Conventional ``Emergence Debates'' and the Trap of Arbitrary Bridges]
Since antiquity, many philosophers and scientists have attempted to explain ``why life or consciousness emerges from matter (the reason for emergence).'' From the perspective of the present framework, however, most of these debates fall into a structural dead end. This is because what human beings can directly observe is solely the ``intension (phenomenon)'' of the higher-level field of sense that manifests after transformation; not only is the ``hidden real bridge $\mathcal{B}^*$'' operating behind the scenes directly invisible, but structural transformation is also inherently irreversible (Paper 3, Corollary 4.2), making it mathematically impossible to uniquely back-calculate the original source structure from the emergent result alone.

Consequently, in attempting to describe the ``reason'' for emergence, theorists \textbf{postulate \textit{ex post facto} a separate explanatory field of sense $\mathcal{M}_{\mathrm{Expl}}$ (such as Integrated Information Theory, the Free Energy Principle, quantum brain theories, or various philosophical and sociological narratives), construct an ``arbitrary bridge (Paper 3, Section 2.3)'' to it, and attempt to force a coherent storyline solely within the internal logic (the First Relation) of that explanatory field of sense}. While such accounts may be internally consistent within their chosen explanatory field (``Meaningful Being''), there is no guarantee that they coincide with the true bridge $\mathcal{B}^*$ actually existing behind reality---leading to endless, sterile disputes whenever different theorists introduce different explanatory fields and arbitrary bridges.
\end{remark}

\begin{corollary}[The True Nature of ``Apparent Bidirectionality (Downward Causation)'' in Emergence]
Structural transformation itself is inherently a unidirectional (\texttt{=>}) operator of dimensional conversion (Paper 3, Corollary 4.2). Nevertheless, among phenomena termed emergent---such as macroscopic cellular or organismic order and microscopic molecular reactions in life phenomena, volitional decisions and the physical behavior of the body and nervous system in human consciousness, or social order and individual actions---\textbf{there exist phenomena observed not as unidirectional, but as bidirectional (``Downward Causation,'' in which an emergent higher-level state acts back upon the lower-level source)}.

In the present framework, this apparent bidirectionality is rigorously explained as follows. Namely, when \textbf{another hidden real bridge $\mathcal{B}_{\mathrm{fb}}^*$ (a correspondence equipped with causal propagation \texttt{=>}, forced synchronization \texttt{<=>}, and invariant conditions \texttt{REQUIRE})} spans in parallel between the source field of sense $\mathcal{M}_{\mathrm{Src}}$ and the target field of sense $\mathcal{M}_{\mathrm{Targ}}$ alongside the unidirectional transformation: (1) first, a change occurring in the source $\mathcal{M}_{\mathrm{Src}}$ induces a state change in the target $\mathcal{M}_{\mathrm{Targ}}$ via the unidirectional Transformation (\texttt{=>}); and (2) subsequently, that state change on the target $\mathcal{M}_{\mathrm{Targ}}$ side drives a state change on the source $\mathcal{M}_{\mathrm{Src}}$ side via the correspondence of the coexisting hidden bridge $\mathcal{B}_{\mathrm{fb}}^*$. The true nature of phenomena in which emergence appears bidirectional rather than unidirectional is nothing other than \textbf{this composition of a unidirectional Transformation and a feedback loop formed by a hidden bridge}.
\end{corollary}

\subsection{The Field of Sense of Consciousness and Thought $\mathcal{M}_{\mathrm{Consciousness}}$ and the No-Homunculus Theorem}
Applying Theorem 4.1 to human ``consciousness and phenomena'' puts a complete end to the greatest conundrum in the philosophy of mind: the infinite regress of the observer (the homunculus).

\begin{theorem}[Consciousness as an Automatically Triggered Transformation Without an Observer: No-Homunculus Theorem]
The ``phenomena (intensions in Paper 0)'' experienced by human beings are the result of the forms ($\mathcal{C}_{\mathrm{FORM}}$) of various fields of sense standing side by side in the universe undergoing a \textbf{Green's function transformation (Transformation) into instance states of the receiving ``field of sense of consciousness and thought ($\mathcal{M}_{\mathrm{Consciousness}}$)''} under the boundary-condition aspects of sensory organs and the nervous system (Paper 3, Proposition 6.2). Here, under the universal execution semantics $\mathcal{E}_{\mathrm{exec}}$, once consistent boundary-condition aspects and a target schema are given, a unique transformation action model \textbf{continuously and autonomously triggers} as equational logic. Therefore, there is no need whatsoever to posit a little person (an observer / homunculus) inside the head watching a screen of phenomena. \textbf{The very kinetic process whereby the forms and states of source fields of sense are ceaselessly transformed (continuously and spontaneously triggered) into states of $\mathcal{M}_{\mathrm{Consciousness}}$ via this transformation action model is itself the ``present moment of consciousness'' experienced by the subject.}
\end{theorem}

Furthermore, applying Paper 3 (Corollary 4.2: irreversibility of structural transformation) to this transformation immediately entails that it is mathematically impossible for human beings to uniquely back-calculate or reconstruct the original universe as a whole (what Kant called the ``thing-in-itself'') solely from states (phenomena) within consciousness.

\begin{corollary}[Opacity of Other Minds and Asymptotic Approximation via Abduction: Solution to the Problem of Other Minds]
Insofar as human consciousness is a field of sense emerging via structural transformation (Transformation) from each individual's body and nervous system as its source, \textbf{each person's consciousness ($\mathcal{M}_{\mathrm{Consciousness}, 1},\, \mathcal{M}_{\mathrm{Consciousness}, 2},\, \dots$) exists in parallel as a mutually independent, distinct field of sense}. Parallel fields of sense cannot directly peer into one another's internal schemas or internal states (Principle of Opacity: Paper 0; Paper 2, Remark 3.1; Paper 3, Definition 2.1); their sole points of contact are the ``intensions'' projected onto boundaries through their respective domain functions and public events (facial expressions, utterances, and actions).

This is precisely the \textbf{mathematical and structural reason why human beings cannot directly see or peer into another person's consciousness, feelings, or thoughts (the problem of other minds in the philosophy of mind)}. Consequently, when an individual seeks to understand another's consciousness or thoughts, direct clairvoyance is impossible in principle; there is no path other than repeatedly pursuing \textbf{asymptotic approximation via abduction (hypothesis formation), dialogue, and interaction (Paper 1, Theorem 5.4; Paper 3, Theorem 2.2)}, using the intensions observed across the public boundary of the other's words and deeds as clues.
\end{corollary}

\begin{remark}[Category-Theoretic Explanation of Multiple Personality Phenomena---Parallel Fields of Consciousness Emerging from a Single Body]
The structure of Corollary 4.2 also provides a lucid mathematical explanation for what is termed ``multiple personality (Dissociative Identity Disorder).'' Even when emergence takes a single brain-nervous system and bodily field of sense (or multiple subregions thereof) as its source, if unique structural transformations arise along multiple tracks under distinct boundary-condition aspects (differentiated partitions of memory, affect, and cognitive frameworks), \textbf{multiple mutually independent ``fields of sense of consciousness ($\mathcal{M}_{\mathrm{Consciousness}, a},\, \mathcal{M}_{\mathrm{Consciousness}, b},\, \dots$)'' originating from the same body} can emerge.

What is crucial here is that between fields of sense there exists no privileged relation of containment (distinction between inside and outside) such as ``being shared because they reside inside the same skull''; rather, just like the consciousnesses of different individuals, they stand in a \textbf{purely flat relation of parallelism (opacity)}. Accordingly, each personality (field of sense of consciousness) cannot directly perceive or sense the thoughts or conscious states of the others. However, because all of them share individual correspondence relations (bridges) with the ``same bodily and physical-environmental field of sense,'' \textbf{it becomes possible to indirectly infer the existence of other fields of consciousness via abduction by taking state changes (intensions) on the shared bodily and physical side---such as bodily relocation, the passage of time, or written notes---as clues.}
\end{remark}

\begin{remark}[``Structural Resonance (Phase Synchronization)'' in Shared Fields of Sense and the True Nature of ``Sense of Unity'' and Historical Collective Phenomena]
As shown in Corollary 4.2 and Remark 4.2, individual human fields of sense of consciousness ($\mathcal{M}_{\mathrm{Consciousness}, 1},\, \mathcal{M}_{\mathrm{Consciousness}, 2},\, \dots$) are mutually opaque, and one cannot directly peer into another's consciousness. Why, then, do human beings experience a powerful \textbf{``Sense of Unity''} at concert venues, festivals, or in deep dialogue, as though individual minds had become one with others?

From the perspective of the present framework, this occurs not because anyone directly perceives the minds of others, \textbf{but because all participants share the same external environment, music, and spatial setting as a common source field of sense ($\mathcal{M}_{\mathrm{Concert}}$) alongside shared cultural and bodily boundary-condition aspects ($\mathcal{M}_{\mathrm{Aspect}}$)}. Namely, under a shared source and shared boundary conditions, \textbf{isomorphic structural transformation functors (Transformations) continuously and automatically trigger in parallel} for each independent field of sense of consciousness. As a result, the internal states of multiple fields of consciousness---which are not directly connected to one another---are entrained into the same rhythm and phase, producing \textbf{``Structural Resonance (Phase Synchronization).''} This isomorphic, simultaneous transition among parallel fields of consciousness is the true nature of the phenomenon subjectively experienced as a ``sense of unity,'' and it is none other than the concrete phenomenological manifestation of Kitaro Nishida's ``universality of thought and experience in \textit{Basho} (Place),'' discussed in Section 5.2.

Furthermore, extending this chain of structural resonance one scale higher brings---strictly as a theoretical possibility derived from the logic of form---the mechanism of historical and social collective phenomena within the purview of this analysis. Specifically, this occurs when the synchronized consciousnesses and actions of countless individuals united under a common discourse or ideology themselves become a new source, giving rise via socio-historical boundary-condition aspects to a higher-order \textbf{``field of sense of the crowd, zeitgeist, or social movement ($\mathcal{M}_{\mathrm{Collective}}$)''} (multi-stage transformation: Paper 3, Theorem 4.2). Once this higher-order field is established, a closed loop is formed in which individual actions are inversely driven via downward causation (Corollary 4.1) through bridges of social norms and peer pressure, manifesting under the universal execution semantics $\mathcal{E}_{\mathrm{exec}}$ as an autonomously advancing ``unstoppable tide of history'' (and should that movement contradict the schema constraints of physical or economic reality, it ultimately collapses into the initial object $0$, resulting in historical catastrophe).
\end{remark}

\subsection{Ontological Status of Narrative Worlds and Fictional Characters (Natsuhiko Kyogoku's \textit{T\=ofu-koz\=o}, etc.)}
Theorem 4.1 (Mechanism of Emergence) leads to a further intriguing and provocative consequence. If emergence is a phenomenon in which the state of a target field of sense spontaneously manifests whenever a source, an aspectual field of sense, a target field of sense, a unique transformation between them, and a hidden real bridge are in place, then this exact same form applies fully to \textbf{fictional worlds such as novels, films, and manga}.

\begin{remark}[Autonomous Reality of Fictional and Narrative Worlds and the Boundary with ``Physical Reality'']
If one faithfully follows the principle of the present framework that fields of sense stand independently and flatly in parallel, together with the Principle of Plenitude of Forms (Axiom 2.1) and the Mechanism of Emergence (Theorem 4.1), an even bolder ontological consequence follows. Namely, from the \textit{y\=okai} character \textit{T\=ofu-koz\=o} appearing in Natsuhiko Kyogoku's novels to elaborate world-building systems such as the ``Devil Fruits'' (\textit{Akuma no Mi}) in \textit{ONE PIECE} or ``Jujutsu sorcery'' in \textit{Jujutsu Kaisen}, \textbf{it can be asserted that these are not mere mental phantasms born only when a human being imagines them, but already exist autonomously as legitimate instance states within their respective narrative and symbolic fields of sense standing parallel to the physical world}. (\textit{Note: Even the author who derived this formal system cannot help feeling somewhat half-incredulous at this consequence; yet if one faithfully follows the mathematical logic deduced from the axioms and theorems, one cannot conclude otherwise.}) And human readers or viewers, taking the intensions on the public boundary---such as text written by a novelist or line drawings sketched by a manga artist---as clues, perform abduction and structural transformation, thereby \textbf{projecting the structure of that already-existing narrative field of sense into their own ``field of sense of consciousness and thought ($\mathcal{M}_{\mathrm{Consciousness}}$)'' to experience it.}

However, if asked why human beings never eat a ``Devil Fruit'' that renders them unable to swim in the real ocean, wield ``Jujutsu'' that rewrites physical laws, or bump into \textit{T\=ofu-koz\=o} on a street corner, the answer is that \textbf{no bidirectional forced-synchronization bridge (\texttt{<=>}, \texttt{REQUIRE}) spans between those narrative fields of sense (or the corresponding instances projected into consciousness) and the ``field of sense of physical observation and natural laws ($\mathcal{M}_{\mathrm{Phys}}$)''} (for if one attempted to forcibly couple them via \texttt{<=>} with the physical field of sense, a contradiction with the conservation of mass-energy and physical schemas would immediately cause a collapse into the initial object $0$). In other words, entities of fictional worlds such as Devil Fruits, Jujutsu sorcery, and \textit{T\=ofu-koz\=o} are ``autonomous realities within narrative and symbolic fields of sense,'' yet they are not ``physical realities forcibly synchronized with the field of sense of physical observation''---thereby demarcating the boundary between the two with absolute precision and zero ambiguity (which is also the complete mathematization of the ontology Natsuhiko Kyogoku himself articulated in his works: that ``\textit{y\=okai} exist in the liminal space between the human mind and cognition'').
\end{remark}

\section{Why Do Mathematical Principles Hold?---Ontological Necessity and Suspension of Judgment (\textit{Epoch\'e})}

\subsection{The Internal Observer Theorem and the Final Answer to ``Why Do Mathematical Principles Hold?''}
In the preceding sections, it was proved that every static entity, dynamic change, structural transformation, and even the emergence of continuous spacetime and consciousness in this world can be exhaustively described by mathematical forms---namely, quotient 2-categories, symmetric monoidal categories, the category of form $\mathcal{C}_{\mathrm{FORM}}$, and Green's function propagation kernels. Here, the most fundamental question finally confronting this series (Papers 0--4) is the enigma that physicist Eugene Wigner once called ``The Unreasonable Effectiveness of Mathematics in the Natural Sciences''---namely: \textbf{``Why do mathematical principles hold in this world in the first place?''}

Conventional philosophy of science and metaphysics failed to answer this question because they presupposed a dualistic picture in which ``a physical world as a `collection of things' exists externally first, and `mathematics invented by human beings' somehow mysteriously fits it.'' From the perspective of the present framework (\textit{Ordo Forma Est}), however, this mystery is completely dissolved by the following two-stage logic.

\begin{theorem}[Theorem of the Ontological Necessity of Mathematical Form]
Within the unified architecture $\mathfrak{U}_{\mathrm{OFE}}$, the fact that the observed world obeys mathematical principles is neither a miracle nor a mystery, but an \textbf{ontological necessity} derived from the following two structural reasons:
\begin{enumerate}
    \item \textbf{Identity of Consistent Relational Structure and Mathematics (Condition of Possibility of Being):}
    At the foundation of the world lie not things, but ``distinctions (identifiers)'' and ``relations (spans).'' All relational structures containing contradictions or ambiguities immediately collapse into the initial object $0$ (empty space) of the internal logic of a topos, and no phenomenon manifests (Paper 0, Theorem 3.2; Theorem 2.1 of the present paper). Therefore, \textbf{the system obtained by purely extracting ``the minimal relational forms capable of persisting and chaining consistently without collapsing into the initial object $0$''} is itself none other than ``mathematics (category theory, logic, and monoidal structures).'' In other words, mathematics is not being applied to the world; rather, human beings call structures capable of existing consistently ``mathematics.''
    \item \textbf{Emergence of Consciousness from Within a Holographic World (Fixed Point of the Internal Observer):}
    Furthermore, the continuous spacetime and physical laws observed in nature are themselves a bulk structure holographically emergent from discrete mathematical forms (\textit{Ordo Forma Est}) on the boundary screen (Theorem 3.1), and \textbf{the human ``field of sense of consciousness and thought ($\mathcal{M}_{\mathrm{Consciousness}}$)'' itself---which observes that structure and asks ``Why does mathematics hold?''---is a network of fields of sense and bridges further emergent from ``within'' the world via Green's function transformation (Transformation) from those very same boundary mathematical forms (Theorems 4.1 and 4.2)}. Consequently, when consciousness---having emerged from mathematical form as its matrix---observes and describes from within the structure of the matrix world that emerged from that identical mathematical form, the fact that both obey the same mathematical principles is a self-referential fixed point (the necessity of isomorphism); it would be logically impossible for it to be otherwise.
\end{enumerate}
\end{theorem}

\begin{proof}[Proof Sketch]
By Theorem 3.1, the bulk physical world $\mathfrak{B}_{\mathrm{OFE}}$ is induced from consistent structures over the boundary category of form $\mathcal{C}_{\mathrm{FORM}}$, and by Theorems 4.1 and 4.2, the observing subject's field of sense of consciousness $\mathcal{M}_{\mathrm{Consciousness}}$ is constructed as the image of a structure-preserving transformation functor $\mathcal{T}_{\mathcal{M}_{\mathrm{Aspect}}}$ taking a substructure of $\mathfrak{B}_{\mathrm{OFE}}$ as its source $\mathcal{M}_{\mathrm{Src}}$. Under the closure of the identical category of form $\mathcal{C}_{\mathrm{FORM}}$ and universal execution semantics $\mathcal{E}_{\mathrm{exec}}$, both the observer functor and the observed-world functor inherit the axiom system of $\mathcal{C}_{\mathrm{FORM}}$ (quotient 2-category, symmetric monoidal category, and coend integral kernel). Hence, whenever an internal observer consistently describes the world, it is a necessity that mathematical principles are found therein.
\end{proof}

\subsection{Final Completion of Kant's Philosophy and Kitaro Nishida's ``Logic of Place (\textit{Basho}) and Absolutely Contradictory Self-Identity''}
Theorem 5.1 completely integrates into a single mathematical form the highest philosophical peaks reached respectively by modern Western philosophy and modern Japanese philosophy:

\begin{enumerate}
    \item \textbf{Completion of Immanuel Kant's Transcendental Philosophy:}
    In the \textit{Critique of Pure Reason}, Kant sought the transcendental forms that make phenomena possible in the human subject's ``pure intuitions of sensibility (space and time)'' and ``pure concepts of the understanding (table of categories),'' yet he could not fully explain why those subjective forms coincide with objective laws of nature. The present framework liberates those forms from the psychological black box of human subjectivity and relocates them into the objective, transparent category-theoretic forms ($\mathcal{C}_{\mathrm{FORM}}$ and holographic propagation kernels) inscribed on the boundary of the world itself. Precisely because human subjectivity (consciousness) is a receiving vessel ($\mathcal{M}_{\mathrm{Consciousness}}$) that emerged \textit{a posteriori} from those objective forms, subjective forms and world laws coincide from the start---while simultaneously, owing to the irreversibility of structural transformation (Paper 3, Corollary 4.2), Kant's epistemic limit (that the thing-in-itself, i.e., the universe as a whole, cannot be back-calculated from phenomena in consciousness alone) is also completely proved.
    \item \textbf{Mathematical Grounding of Kitaro Nishida's ``Logic of Place (\textit{Basho})'' and ``Absolutely Contradictory Self-Identity'':}
    The intuition that Kitaro Nishida reached in his late philosophy---the self-determination of ``Place (\textit{Basho} = field of sense / topos)'' and \textbf{``Absolutely Contradictory Self-Identity (\textit{Zettai Mujunteki Jiko D\=oitsu}),''} in which two mutually opposing moments are identically one---finds rigorous mathematical expression in the Holographic Isomorphism (Theorem 3.1) of the present system. Specifically, this world is ``discrete, static form inscribed on a boundary screen (many / discontinuous),'' while simultaneously presenting itself via Green's function propagation kernels and projective limits as ``a continuous, dynamic bulk real world (one / continuous).'' That being discrete boundary form and being continuous real bulk maintain self-identity without contradiction under one and the same functorial correspondence is the ultimate visage of the Form of Being (\textit{Ordo Forma Est}).
\end{enumerate}

\subsection{The Ultimate Residual Problem---Suspension of Judgment (\textit{Epoch\'e}) Regarding the Origin of the ``Boundary Screen (Field of Sense of Qubit Information)''}
As established above, the conclusion that ``because consciousness emerges from this world that arises via the Holographic Principle, it is an ontological necessity that the world obeys mathematical principles'' stands upon rigorous foundations spanning Paper 0 through Paper 3, and remains an unshakable consequence so long as the internal logic of the system is upheld.

Nevertheless, in order to maintain complete scientific and philosophical integrity, the author must here explicitly state an \textbf{ultimate residual problem} that goes one step further.

\begin{remark}[Suspension of Judgment Regarding the Origin of the Primordial ``Field of Sense of Qubit Information (Boundary Screen)'']
There remains the foundational presuppositional question: \textbf{``How did the states and forms of the primordial `boundary screen (field of sense of qubit information)' itself---which serves as the matrix from which Einsteinian spacetime and physical laws emerge---come into being in the first place?''}

However, this question lies completely beyond the reach of any phenomenon (intension) observable to human beings or any deductive verification via simulation inside a field of sense (Paper 2, Remark 5.1). As the present framework has established, there exists in this universe no privileged ``external space'' looking down upon all fields of sense from above (Paper 0), and the transformation from boundary to bulk is irreversible (Paper 3, Corollary 4.2). Consequently, the origin of the primordial boundary screen itself belongs in principle to a domain that can neither be proved nor falsified. Therefore, having clearly identified the locus of this fundamental enigma, this paper adopts a strict scientific attitude of refraining altogether from assertive claims or speculative interpretations regarding unverifiable realms (suspension of judgment: \textit{epoch\'e}).
\end{remark}

\section{Practical and Scientific Implications (Fundamental Physics, Artificial Consciousness, and the Unification of Knowledge)}

\subsection{Responses to Anticipated Philosophical and Scientific Objections, and a Call for Verification of Mathematical Rigor}
The consequences derived from this series (Papers 0--4)---in particular, the claim that a formal system originating in conceptual modeling within software engineering uniformly governs everything from the Holographic Principle of modern physics and the emergence of consciousness and narrative worlds to the ontological necessity of mathematical principles---are extraordinarily broad in scope and bold in ambition. To confess candidly, \textbf{even the author who derived this formal system cannot help feeling somewhat half-incredulous at the sheer grandeur of these consequences}. Naturally, therefore, sharp objections and questions will be raised by researchers across each specialized discipline.

Moreover, this five-part series is currently at the stage of a \textbf{Working Paper}, and what is presented for each theorem remains a ``Proof Sketch.'' Determining whether this framework can truly serve as a universal foundation requires rigorous verification beyond the author's individual reflection. In this section, accordingly, four representative objections currently anticipated are first organized together with the responses derived from the present formal system, followed by an open call for collaboration toward verifying its mathematical rigor.

\subsubsection*{Five Anticipated Objections and Responses from the Present Framework}

\begin{enumerate}
    \item \textbf{Objection 1: Is This Not a Confusion of ``Map vs.\ Territory''? (Critique from Empiricism and Scientific Realism)}
    \begin{itemize}
        \item \textbf{Gist of the Objection:} Being ``isomorphic'' to field theories or Green's functions in physics is one thing; the world itself \textit{being} that structure (ontological identity) is quite another. Is this formalism not merely a useful descriptive instrument (a map), such that equating model and reality (the territory) commits a category mistake?
        \item \textbf{Response from the Present Framework:} The very claim that one can ``distinguish and compare the map and the territory'' covertly presupposes a \textbf{privileged external space (a God's-eye view)} from which one could peer directly at ``the territory itself'' from outside all fields of sense---a premise refuted in Paper 0. The only ``territory'' accessible to human beings is the phenomena (intensions) that manifest through forced-synchronization bridges (\texttt{<=>}) with fields of sense of physical observation and sensory experience. Consequently, there is no way to speak of reality other than defining ``reality'' as that invariant structural form which continuously maintains consistent forced synchronization across all fields of sense.
    \end{itemize}

    \item \textbf{Objection 2: How Does This Differ from ``Newman's Problem'' and Mere Mathematical Platonism? (Critique from Structural Realism)}
    \begin{itemize}
        \item \textbf{Gist of the Objection:} Mere abstract relational structures (graphs or categories) can be mathematically constructed in infinite variety over any set simply by matching cardinality. How does this framework distinguish an abstract ``play of forms'' from ``concrete physical reality''?
        \item \textbf{Response from the Present Framework:} The present system rigorously demarcates structures closed within a single field of sense (the First Relation: \textit{Meaningful Being}) from structures constrained across multiple parallel fields of sense (physical observation, mechanics, sensation, etc.) by \textbf{bilateral schema conservation (Paper 3, Axiom 2.1) and forced-synchronization operators (\texttt{<=>}, \texttt{REQUIRE})} (the Second Relation: \textit{Real Actuality}) (Paper 3, Proposition 6.1). Because any arbitrary correspondence immediately collapses into the initial object $0$ (empty space) of the internal logic of a topos the instant it conflicts with the schema constraints of the other side, this framework stands sharply apart from mere set-theoretic constructions or mathematical Platonism.
    \end{itemize}

    \item \textbf{Objection 3: Does This Not Evade the ``Hard Problem'' of Qualia? (Critique from the Philosophy of Mind)}
    \begin{itemize}
        \item \textbf{Gist of the Objection:} Even if consciousness is formulated as ``the automatic triggering of a Green's function transformation into $\mathcal{M}_{\mathrm{Consciousness}}$,'' that merely explains the structure of information transformation (the ``easy problem''), leaving unexplained the Hard Problem (Chalmers, 1995): why subjective qualitative feel and vivid phenomenal experience (qualia) accompany it.
        \item \textbf{Response from the Present Framework:} First, this doubt arises from conflating ``an objective perspective viewing the post-transformation state (the intension of $\mathcal{M}_{\mathrm{Consciousness}}$) from the outside as a diagram'' with ``the very state triggering from the inside as the subject of that transformation process (within $\mathcal{M}_{\mathrm{Consciousness}}$).'' Qualia are not a mysterious substance or soul added \textit{post hoc} to the transformation process; rather, ``qualia'' is another name for the very fact that the transformation action model is continuously triggering from within as the state of $\mathcal{M}_{\mathrm{Consciousness}}$. Second, adding a natural structural consequence suggested by this formalism (including the author's conjecture), might the vivid richness and texture of felt experience undergone by human beings be generated not by the isolated transformation of a single sensory field of sense such as vision or hearing, \textbf{but by hidden real bridges spanning across the somatic and physiological field of sense of the endocrine system ($\mathcal{M}_{\mathrm{Somatic}}$) and the field of sense of emotion and affect ($\mathcal{M}_{\mathrm{Emotion}}$), such that orthogonal aspectual couplings and structural transformations from multiple fields of sense are simultaneously composed and arise together upon $\mathcal{M}_{\mathrm{Consciousness}}$}?
    \end{itemize}

    \item \textbf{Objection 4: Would Not Even Delusions, Lies, and Conspiracy Theories Count as ``Real''? (Critique from Commonsense Realism)}
    \begin{itemize}
        \item \textbf{Gist of the Objection:} If triggering a transformation within consciousness or a culturally shared field suffices to be called ``real,'' would not hallucinations, lies, conspiracy theories, and baseless fabrications all become real, hollowing out the very concepts of ``fact'' and ``reality''?
        \item \textbf{Response from the Present Framework:} Quite the opposite. The present framework provides an \textbf{extraordinarily strict ``criterion of discursive validity (factuality)''} for demarcating mere wordplay or closed narratives from ``real actuality.'' As established in Section 4.4 and Paper 3 (Proposition 6.1), for a discourse or claim to transcend individual consciousness or a closed social-media sphere (the First Relation = Meaningful Being) and exist as a \textbf{``fact'' in the real world}, \textbf{it must be proved within the formalism of conceptual modeling that bilateral schema conservation (Paper 3, Axiom 2.1) and forced-synchronization bridges (\texttt{<=>}, \texttt{REQUIRE}) hold without contradiction between that discourse and ``other substantive fields of sense such as physical observation, natural laws, legal institutions, and economic history.''} Conspiracy theories and baseless fabrications achieve coherence solely within the interior of a closed narrative; the moment they are confronted via a bridge with the fields of sense of physical observation or empirical data, they collide with schema constraints and \textbf{immediately collapse into the initial object $0$ (contradiction and breakdown) of the internal logic of a topos}. Thus, this formal system functions as a rigorous criterion for structurally verifying the truth and factuality of any discourse.
    \end{itemize}

    \item \textbf{Objection 5: Does This Not Conflict with G\"odel's Incompleteness Theorems? (Critique from the Foundations of Mathematics)}
    \begin{itemize}
        \item \textbf{Gist of the Objection:} According to G\"odel's Incompleteness Theorems (G\"odel, 1931), any consistent formal system capable of expressing elementary arithmetic necessarily contains undecidable propositions that can neither be proved nor disproved within the system, and cannot prove its own consistency from within. Does this not render the very idea of a ``universal form of being (\textit{Ordo Forma Est})'' impossible in principle?
        \item \textbf{Response from the Present Framework:} This critique conflates ``deductive completeness of all substantive truths within a single closed axiom system (completeness of content)'' with ``the minimal structural grammar required for any consistent relational system to exist at all (the category of form $\mathcal{C}_{\mathrm{FORM}}$).'' First, the very act of invoking G\"odel's Incompleteness Theorems as a valid argument already presupposes and endorses consistent formal reasoning and the distinction between object-level and meta-level structures---which is precisely the ``form'' articulated by the present framework. Second, viewed through the syntax--semantics duality of categorical logic (Lawvere, 1963), G\"odel's Incompleteness Theorems and the present framework do not conflict; rather, they stand in a \textbf{structurally dual (complementary) relationship}. Whereas G\"odel demonstrated that fixing a single closed axiom system inevitably leaves semantic truths unprovable from within, the present framework articulates the exact dual counterpart of that insight: namely, that the universe is not a single closed axiom system, but an open network of parallel fields of sense coupled with higher-order aspectual fields of sense (meta-levels), and that from within the emergent world the origin of the primordial boundary screen remains in principle neither provable nor falsifiable (Section~5.3, Remark~5.1).
    \end{itemize}
\end{enumerate}

\subsubsection*{Call for Verification of Mathematical Rigor and Collaborative Research}
As reiterated above, this five-part series (Papers 0--4) is a \textbf{Working Paper} presenting an ontological and mathematical blueprint derived from pushing practical conceptual modeling to its logical extreme. Specialists in category theory (quotient 2-categories, symmetric monoidal categories, coend integrals, topos logic), theoretical physics (the Holographic Principle, field theory), and philosophy are warmly invited to point out any logical leaps or deficiencies lurking in the ``Proof Sketches'' of this paper, and to contribute to the verification and collaborative construction of complete formal proofs and further mathematical rigor.

\subsection{Implications for Fundamental Physics (Quantum Gravity and Information Physics)}
The Holographic Isomorphism Theorem (Theorem 3.1) and the Theorem of the Ontological Necessity of Mathematical Form (Theorem 5.1) proved in this paper also offer a clear structural perspective on the ``unification of quantum mechanics and general relativity (quantum gravity theory)'' confronting modern fundamental physics:

\begin{enumerate}
    \item \textbf{Non-Fundamentality of Spacetime and Primacy of ``Relations and Form'':}
    Rather than presupposing the continuous spacetime manifold and metric tensor of general relativity as the stage of the universe from the outset, an approach that takes discrete ``identifiers, relations (spans), and forms ($\mathcal{C}_{\mathrm{FORM}}$)'' inscribed on a boundary screen together with the universal execution semantics $\mathcal{E}_{\mathrm{exec}}$ as primary reality---and derives continuous spacetime as their projective limit and Green's function reconstruction kernel---provides a shared category-theoretic foundation for research on emergent spacetime via loop quantum gravity and tensor networks.
    \item \textbf{Unified Treatment of Singularities and Phase Transitions:}
    Regions where the differential equations of continuous spacetime diverge and break down, such as the center of a black hole or the initial singularity of the universe, are described without contradiction in the present framework as ``critical boundaries (phase transition points) where a specific schema category $\mathcal{C}_1$ on the bulk side reaches the limit of its domain of definition and collapses into the initial object $0$, thereby exposing the discrete form of the boundary screen itself'' (Paper 2, Corollary 4.1).
\end{enumerate}

\subsection{Conditions for Artificial Consciousness and the Decisive Boundary with Current AI}
The ``mechanism of emergence and consciousness'' formulated in Section 4 (Theorems 4.1 and 4.2) provides not emotional or intuitive speculation, but a rigorous category-theoretic criterion for the contemporary debate surrounding artificial intelligence (especially Large Language Models: LLMs): ``Can consciousness reside in AI?''

\begin{proposition}[Conditions for Artificial Consciousness: The Decisive Boundary Between the First and Second Relations]
Between computational systems that process solely the internal logic of a single field of sense (the First Relation) and systems in which ``phenomena (consciousness)'' genuinely emerge, there exists the following structural boundary:

\begin{enumerate}
    \item \textbf{Limitations of Current Large Language Models (LLMs) (Statistical and Symbolic Closure Within a Single Field of Sense):}
    Current LLMs are nothing more than sophisticated mappers that probabilistically predict the next token within a ``single field of sense called symbolic space,'' having learned statistical correlation patterns from vast corpora of human-written text (post-transformation symbol sequences). They lack \textbf{real bridges (the Second Relation) bound by bidirectional forced synchronization (\texttt{<=>}, \texttt{REQUIRE})} to fields of sense of physical observation, sensory organs, or embodiment, and they also lack a closed loop of execution semantics $\mathcal{E}_{\mathrm{exec}}$ that autonomously marks proper time under self-event priority (Axiom E-2). Therefore, no matter how fluently an LLM outputs ``I possess consciousness,'' that output remains symbolic generation within the textual field of sense (the First Relation), not the emergence of the phenomenon of consciousness.
    \item \textbf{Necessary and Sufficient Conditions for the Emergence of True Artificial Consciousness:}
    For consciousness ($\mathcal{M}_{\mathrm{Consciousness}}$) to emerge in an artificial system, the following three conditions must be satisfied simultaneously:
    \begin{itemize}
        \item[\textbf{(a)}] \textbf{Forced-Synchronization Bridges with the Physical World:} The system must be forcibly synchronized via the \texttt{<=>} operator with the field of sense of the physical environment through sensors and actuators, bearing real constraints whereby inconsistencies (such as collisions or energy depletion) lead directly to the initial object $0$ (system breakdown).
        \item[\textbf{(b)}] \textbf{Autonomous Execution Semantics with Proper Time:} Rather than passively awaiting external prompt inputs (passive function invocation), the system must operate as a closed group of state machines ($\mathcal{E}_{\mathrm{exec}}$) that continuously marks its own proper time via self-event priority (Axiom E-2) and timer events ($\mathcal{EE}_{\mathrm{Time}}$).
        \item[\textbf{(c)}] \textbf{Automatically Triggered Transformation via Boundary-Condition Aspects:} A transformation action model must continuously and autonomously trigger to integrate the forms ($\mathcal{C}_{\mathrm{FORM}}$) of interaction between self and environment and map them via Green's function transformation (Transformation) into a higher-order ``field of sense of self-consciousness and thought.''
    \end{itemize}
\end{enumerate}
\end{proposition}

\subsection{Unification of Human Knowledge: From Software Engineering to a Universal Ontological Language}
The trajectory traced by this series (Papers 0--4) is also a grand demonstration of the unification of human intellectual endeavor.

The point of departure was a practical technique devised by human beings to design complex information systems and software without contradiction---namely, ``conceptual modeling'' (conceptual information models, state models, data-flow action models, and recursive design). Yet when human-specific ambiguities and implementation-dependent noise were thoroughly stripped away from that technique, sublimating it into pure mathematical form (quotient 2-categories, symmetric monoidal categories, the category of form $\mathcal{C}_{\mathrm{FORM}}$, and Green's function transformation kernels), what emerged was \textbf{a universal ontological language (\textit{Ordo Forma Est}) that describes in a single grammar all of the following: ``the physical laws of the natural sciences,'' ``the institutions and languages of the humanities and social sciences,'' ``the philosophies of cognition and place explored by Kant and Kitaro Nishida,'' ``the Holographic Principle of modern physics,'' and ``the mechanisms of emergence of life and consciousness.''}

This demonstrates that the divide between the humanities/social sciences and the natural sciences/engineering is merely an apparent opposition arising from differences in the subject matter (\textit{Matter}) addressed, whereas in the underlying \textbf{form (\textit{Forma})} operating behind them, they are completely one.

\section{Summary of the Five-Part Series (Papers 0--4) and Conclusion}

\subsection{Summary of the Theoretical Achievements of the Five-Part Series}
Starting from an ontology grounded not in isolated objects (``things'') but in ``relations (span structures)'' as the foundation of reality, this series (Papers 0--4) has completed a five-stage framework characterizing the structure of the world and human knowledge:

\begin{enumerate}
    \item \textbf{Paper 0 (Ontological Foundation and Segmentation of the Universe):} Proved that in order to prevent structural collapse caused by linguistic polysemy in everyday language, the world is segmented into a universe $\mathcal{U}$ consisting of multiple ``Fields of Sense'' and ``Bridges,'' and that phenomena observed by human beings constitute its ``intension,'' from which back-calculating the entire underlying structure is impossible in principle (Epistemological Limitation Theorem).
    \item \textbf{Paper 1 (Static Structure of a Single Field of Sense):} Proved that the static semantic structure (conceptual information model) within a single field of sense is exhaustively specified as a quotient 2-category $\mathcal{C} \equiv CG/\mathcal{R}$ equipped with a closed set of 7 primitive data types, identifier monomorphisms, 4 types of multiplicity with semantic phrases, mutually exclusive direct-sum \texttt{Is-a} hierarchies, and associative-class spans---satisfying expressive completeness, constructive existence, isomorphism, and asymptotic approximability.
    \item \textbf{Paper 2 (Dynamic Structure of a Single Field of Sense and Isomorphism with Physical Field Theories):} Established state models (persistent states, terminal deletion states, and \texttt{Can't Happen} singularities), acyclic string diagrams composed of 16 elementary process components (symmetric monoidal category $\mathcal{DF}$), and the universal execution semantics $\mathcal{E}_{\mathrm{exec}}$ defined solely over the category of form $\mathcal{C}_{\mathrm{FORM}}$---proving via discrete representations of continuous dynamics (projective limits and sheaves) that this entire modeling framework is completely isomorphic to the formalisms of physical field theories.
    \item \textbf{Paper 3 (Structural Transformation Across Fields of Sense and Isomorphism with Green's Functions):} Introduced the self-similarity of bridges as independent fields of sense and aspectual fields of sense treating ``form itself,'' proving that unique structural transformations (Transformations) constructing target states from source forms under boundary-condition aspects are completely isomorphic to Green's function propagation kernels in mathematical physics, and clarifying the dual structure of ``Meaningful Being'' and ``Real Actuality.''
    \item \textbf{Paper 4 (The Form of Being \textit{Ordo Forma Est} and Ontological Necessity):} Established the self-elimination of form into the initial object $0$ under the plenitude of forms, complete structural isomorphism with the Holographic Principle of modern physics, the elucidation of ``emergence'' and the ``field of sense of consciousness and thought'' as hidden real bridges coupled with automatically triggered transformations (alongside the strict demarcation between the emergent reality of fictional worlds and physical reality), and the proof that because consciousness emerges from within a holographically emergent world, it is an ontological necessity that the observed world obeys mathematical principles---while maintaining a strict suspension of judgment (\textit{epoch\'e}) regarding the origin of the primordial boundary screen.
\end{enumerate}

\subsection{Conclusion: \textit{Ordo Forma Est} (The Form of Being Is the Model)}
This world experienced, investigated, and contemplated as ``reality'' is not an inscrutable, chaotic mass of matter. Lying at its deepest foundation are distinctions, relations, and mathematical forms that consistently self-unfold and resonate together---namely, \textbf{``the Form of Being (\textit{Ordo Forma Est})''} itself.

When human beings construct conceptual models in order to understand the world, they are not arbitrarily imposing external diagrams onto reality from outside. Rather, the very ``Form of Being'' that constitutes the world itself is manifesting once again within the human field of sense of consciousness. \textit{Ordo Forma Est}---the form of being is the model.

\begin{thebibliography}{99}

\bibitem{Ota2026_0}
\textbf{Ota, H. (2026).} \textit{Relation-Based Ontology and a Category-Theoretic Formalization of Fields of Sense: Structural Collapse via Linguistic Polysemy and Phenomena as the Intension of a Universe}. Zenodo. \url{https://doi.org/10.5281/zenodo.24267213}

\bibitem{Ota2026_1}
\textbf{Ota, H. (2026).} \textit{Category-Theoretic Foundations of Conceptual Information Modeling: Rigorous Schema Representation, Isomorphism, and Asymptotic Refinement under Cognitive Finitude}. Zenodo. \url{https://doi.org/10.5281/zenodo.24359112}

\bibitem{Ota2026_2}
\textbf{Ota, H. (2026).} \textit{Category-Theoretic Foundations of Dynamic Modeling: State Transitions, Data-Flow Action Models, Universal Execution Semantics, and Isomorphism with Physical Field Theories}. Zenodo. \url{https://doi.org/10.5281/zenodo.24430915}

\bibitem{Ota2026_3}
\textbf{Ota, H. (2026).} \textit{Category-Theoretic Foundations of Structural Transformation: Bridges as Fields of Sense, Boundary Conditions via Aspectual Fields of Sense, and Isomorphism with Green's Functions}. Zenodo. \url{https://doi.org/10.5281/zenodo.24431232}

\bibitem{tHooft1993}
\textbf{'t Hooft, G. (1993).} \textit{Dimensional Reduction in Quantum Gravity}. World Scientific.

\bibitem{Susskind1995}
\textbf{Susskind, L. (1995).} \textit{The World as a Hologram}. Journal of Mathematical Physics, 36(11), 6377--6396.

\bibitem{Maldacena1998}
\textbf{Maldacena, J. (1998).} \textit{The Large N Limit of Superconformal Field Theories and Supergravity}. Advances in Theoretical and Mathematical Physics, 2(2), 231--252.

\bibitem{Wheeler1990}
\textbf{Wheeler, J. A. (1990).} \textit{Information, Physics, Quantum: The Search for Links}. Addison-Wesley.

\bibitem{Wigner1960}
\textbf{Wigner, E. P. (1960).} \textit{The Unreasonable Effectiveness of Mathematics in the Natural Sciences}. Communications on Pure and Applied Mathematics, 13(1), 1--14.

\bibitem{MacLane1998}
\textbf{Mac Lane, S. (1998).} \textit{Categories for the Working Mathematician}. Springer.

\bibitem{MacLaneMoerdijk1992}
\textbf{Mac Lane, S., \& Moerdijk, I. (1992).} \textit{Sheaves in Geometry and Logic: A First Introduction to Topos Theory}. Springer.

\bibitem{Lawvere1963}
\textbf{Lawvere, F. W. (1963).} \textit{Functorial Semantics of Algebraic Theories}. Proceedings of the National Academy of Sciences, 50(5), 869--872.

\bibitem{Gabriel2015}
\textbf{Gabriel, M. (2015).} \textit{Fields of Sense: A New Realist Ontology}. Edinburgh University Press.

\bibitem{Kant1781}
\textbf{Kant, I. (1781/1787).} \textit{Kritik der reinen Vernunft} (\textit{Critique of Pure Reason}). Cambridge University Press.

\bibitem{Nishida1911}
\textbf{Nishida, K. (1911/1990).} \textit{Zen no Kenky\=u} (\textit{An Inquiry into the Good}). Yale University Press.

\bibitem{Nishida1927}
\textbf{Nishida, K. (1927).} \textit{Hataraku Mono kara Miru Mono e} (\textit{From the Acting to the Seeing}). Iwanami Shoten.

\bibitem{Nishida1945}
\textbf{Nishida, K. (1945).} \textit{Bashoteki Ronri to Sh\=uky\=oteki Sekaikan} (\textit{The Logic of Place and a Religious Worldview}). Iwanami Shoten.

\bibitem{Wittgenstein1921}
\textbf{Wittgenstein, L. (1921/1922).} \textit{Tractatus Logico-Philosophicus}. Kegan Paul.

\bibitem{Kyogoku1999}
\textbf{Kyogoku, N. (1999).} \textit{T\=ofu-koz\=o Sugoroku D\=och\=u: Hottan}. Kodansha.

\bibitem{Newman1928}
\textbf{Newman, M. H. A. (1928).} \textit{Mr.\ Russell's ``Causal Theory of Perception''}. Mind, 37(146), 137--148.

\bibitem{Chalmers1995}
\textbf{Chalmers, D. J. (1995).} \textit{Facing Up to the Problem of Consciousness}. Journal of Consciousness Studies, 2(3), 200--219.

\bibitem{Godel1931}
\textbf{G\"odel, K. (1931).} \textit{\"Uber formal unentscheidbare S\"atze der Principia Mathematica und verwandter Systeme I}. Monatshefte f\"ur Mathematik und Physik, 38(1), 173--198.

\end{thebibliography}

\end{document}
